\chapter{Herramientas, compilación y pruebas}
\label{cap:herramientas}

\section{Lo que hace falta}

\begin{center}
\begin{tabularx}{\linewidth}{@{}lXl@{}}
\toprule
Pieza & Para qué & Cuándo \\
\midrule
Rust (cargo) & compilar el núcleo & al compilar \\
Compilador de C++ & el puente y el código nativo & al compilar \\
Qt 6.8+ (base, declarative, shadertools) & la interfaz & al compilar y al usar \\
mpv & reproducir & al usar \\
yt-dlp + deno (o Node.js) & obtener audio cuando la vía rápida no basta & al usar \\
qt6-webengine (solo Linux, opcional) & ventana de inicio de sesión & al compilar y al usar \\
\bottomrule
\end{tabularx}
\end{center}

Los comandos exactos para macOS y Arch están en \archivo{COMPILAR.md}.

\section{\texttt{Cargo.toml}: el paquete}

\begin{codigo}
[package]
name = "musicplayer"
edition = "2021"

[workspace]
members = ["login"]
default-members = [".", "login"]

[dependencies]
cxx = "1.0.95"
cxx-qt = "0.10"
cxx-qt-lib = { version = "0.10", features = ["qt_gui", "qt_qml"] }
dirs = "6"
md5 = "0.8"
serde = { version = "1", features = ["derive"] }
serde_json = "1"
sha1 = "0.10"
souvlaki = { version = "0.9", default-features = false, features = ["use_zbus"] }
ureq = { version = "3", features = ["json"] }

[build-dependencies]
cc = "1"
cxx-qt-build = { version = "0.10", features = ["link_qt_object_files"] }

[profile.release]
opt-level = 3
lto = "thin"
strip = true
\end{codigo}

\subsection{Las dependencias, una por una}

Cada una está por un motivo; el proyecto evita dependencias que no se ganen su sitio.

\begin{description}[leftmargin=1.2em]
  \item[\cod{cxx}, \cod{cxx-qt}, \cod{cxx-qt-lib}] El puente con Qt
  (capítulo~\ref{cap:puente}). \cod{cxx-qt-lib} trae los tipos de Qt (\cod{QString},
  \cod{QGuiApplication}, \cod{QQmlApplicationEngine}…).
  \item[\cod{serde}, \cod{serde_json}] Convertir structs a JSON y desde JSON. Se usan para hablar
  con YouTube, con mpv, con Discord y con la interfaz.
  \item[\cod{ureq}] Cliente HTTP \emph{síncrono} y pequeño. Síncrono a propósito: el proyecto
  usa hilos normales en vez de programación asíncrona (\cod{async}/\cod{tokio}), que añadiría un
  motor entero y bastante complejidad para un número de peticiones muy pequeño.
  \item[\cod{dirs}] Rutas estándar de cada sistema: configuración
  (\texttt{\textasciitilde/Library/Application Support} en macOS, \texttt{\textasciitilde/.config}
  en Linux) y caché.
  \item[\cod{sha1}] La firma de las peticiones con sesión a YouTube (SAPISIDHASH).
  \item[\cod{md5}] La firma de las peticiones a Last.fm (su API la exige).
  \item[\cod{souvlaki}] Controles multimedia del sistema: el «Ahora suena» del Centro de
  control de macOS y MPRIS en Linux (las teclas multimedia, el widget del escritorio). Se
  compila con \cod{use_zbus} para no depender de las bibliotecas de D-Bus en C.
  \item[\cod{cc}] (al compilar) Compilar el C++ y el Objective-C++ propios.
  \item[\cod{cxx-qt-build}] (al compilar) Generar el código del puente y enlazar con Qt.
\end{description}

\subsection{El perfil \texttt{release}}

\cod{opt-level = 3} optimiza al máximo; \cod{lto = "thin"} permite optimizar entre crates (algo
más lento de compilar, binario más rápido y pequeño); \cod{strip = true} quita los símbolos de
depuración del binario (pesa menos y carga antes).

\subsection{El espacio de trabajo}

El repositorio tiene dos paquetes: el reproductor (la raíz) y \archivo{login/}
(\cod{musicplayer-login}, la ventana de inicio de sesión de Linux). \cod{[workspace]} los
compila juntos con un solo \cod{cargo build}; \cod{default-members} hace que se compilen los
dos por defecto.

\begin{porque}[¿Por qué un binario aparte para el inicio de sesión de Linux?]
La ventana de Linux usa Qt WebEngine, que es Chromium entero. Si el reproductor lo enlazara,
lo cargaría al arrancar aunque nunca se use. En un binario aparte, solo ocupa memoria
mientras la ventana está abierta. Además, si falta \texttt{qt6-webengine} el reproductor se
compila igual.
\end{porque}

\section{\texttt{build.rs}: lo que pasa antes de compilar}

Un \archivo{build.rs} es un programa en Rust que cargo ejecuta \emph{antes} de compilar el
paquete. El del reproductor hace tres cosas:

\begin{rust}
const QML_FILES: &[&str] = &["qml/Main.qml", "qml/Ambient.qml", /* ... */];
const QML_SINGLETONS: &[&str] = &["qml/Theme.qml", "qml/Catalog.qml"];

fn main() {
    // 1) Un módulo QML «app.musicplayer» con todos los archivos QML.
    let module = QmlModule::new("app.musicplayer")
        .qml_files(QML_FILES.iter().copied())
        .qml_files(QML_SINGLETONS.iter().map(|path| QmlFile::from(path).singleton(true)));

    // 2) El puente: genera C++ para cada archivo Rust con #[cxx_qt::bridge] y enlaza Qt.
    CxxQtBuilder::new_qml_module(module)
        .qt_module("Gui").qt_module("Quick").qt_module("Network")
        .cpp_file("src/window.cpp")
        .files(["src/account.rs", "src/library.rs", "src/player.rs",
                "src/plugins.rs", "src/settings.rs"])
        .build();

    // 3) Solo en macOS: la ventana de inicio de sesión en Objective-C++.
    if std::env::var("CARGO_CFG_TARGET_OS").as_deref() == Ok("macos") {
        cc::Build::new().cpp(true).file("src/login_mac.mm").flag("-fobjc-arc")
            .flag("-std=c++17").compile("musicplayer_login");
        println!("cargo::rustc-link-lib=framework=Cocoa");
        println!("cargo::rustc-link-lib=framework=WebKit");
    }
}
\end{rust}

\begin{itemize}
  \item Los archivos QML se \textbf{compilan dentro del binario} (como recursos
  \texttt{qrc:/}). No hay que copiar nada junto al ejecutable, y QML los precompila, así que
  arranca antes.
  \item El módulo se llama \cod{app.musicplayer}: es lo que importan los QML
  (\cod{import app.musicplayer}).
  \item Las líneas \cod{println!("cargo::...")} son instrucciones para cargo: enlazar una
  biblioteca, recompilar si cambia un archivo o una variable de entorno, etc.
\end{itemize}

\begin{trampa}
Un archivo QML nuevo que no se añade a \cod{QML_FILES} \emph{no existe} para QML («Type X
unavailable»). Un archivo Rust nuevo con \cod{#[cxx_qt::bridge]} hay que añadirlo a
\cod{.files([...])} y como \cod{mod} en \archivo{src/main.rs}.
\end{trampa}

\section{Compilar y ejecutar}

\begin{terminal}
cargo build --release          # compila (la primera vez, unos minutos; luego ~15 s)
./target/release/musicplayer   # ejecuta
cargo run --release            # las dos cosas
\end{terminal}

Compila siempre en \texttt{release} si vas a medir velocidad o memoria: la versión de
depuración (\texttt{cargo build}) es varias veces más lenta.

Para tenerla como aplicación:

\begin{description}[leftmargin=1.2em]
  \item[macOS: \archivo{scripts/crear-app-macos.sh}] Compila, crea \texttt{Reproductor.app}
  con su icono e \texttt{Info.plist}, la firma con una firma local (Apple Silicon no abre apps
  sin firmar) y la copia a \texttt{/Applications}. Con \texttt{--independiente} mete Qt dentro
  de la app (\texttt{macdeployqt}), para que no dependa de Homebrew.
  \item[Linux: \archivo{scripts/instalar-linux.sh}] Compila e instala para tu usuario en
  \texttt{\textasciitilde/.local} (programa, entrada de menú, icono y, si se compiló con
  WebEngine, \texttt{musicplayer-login}). Con \texttt{--dependencias} instala antes los paquetes
  con pacman.
\end{description}

\section{Probar sin abrir la ventana: las sondas}

\archivo{src/main.rs} tiene un modo «sonda»: hace una consulta, imprime el resultado y sale,
sin crear la interfaz. Es la forma más rápida de saber si un problema es de la red o de la
interfaz.

\begin{center}
\small
\begin{longtable}{@{}p{0.42\linewidth}p{0.52\linewidth}@{}}
\toprule
Orden & Qué hace \\
\midrule
\endhead
\texttt{--probe home} / \texttt{explore} / \texttt{albums} & Secciones de YouTube Music \\
\texttt{--probe playlists} & Playlists por estado de ánimo \\
\texttt{--probe search "texto"} & Búsqueda \\
\texttt{--probe queue ID} & Radio de una canción \\
\texttt{--probe album ID} & Canciones de un álbum o playlist (página real) \\
\texttt{--probe lyrics "título" "artista"} & Letra sincronizada \\
\texttt{--probe play ID} & Resuelve y reproduce 25 s, con tiempos \\
\texttt{--probe resolve id1,id2} & Mide la resolución (con el yt-dlp residente) \\
\texttt{--probe fast id1,id2} & Solo la vía rápida, y por qué falla cada una \\
\texttt{--probe url ID} & La URL de la vía rápida \\
\texttt{--probe playable ID} & Qué dice el reproductor de YouTube Music de una pista \\
\texttt{--probe account} & Si YouTube reconoce la sesión guardada \\
\texttt{--probe cookies chrome} & Diagnóstico de importar la sesión de un navegador \\
\bottomrule
\end{longtable}
\end{center}

\section{Capturas automáticas}

\archivo{qml/Main.qml} acepta argumentos de depuración. El principal es \texttt{--shot}: abre
la ventana, espera, guarda una captura, escribe una línea \texttt{estado:} con el estado del
reproductor y sale.

\begin{terminal}
MUSICPLAYER_SILENT=1 ./target/release/musicplayer --shot /tmp/c.png --delay 12000 \
  --section albumes --index 3 --autoplay --autoplay-delay 4000 --now
\end{terminal}

\begin{center}
\small
\begin{longtable}{@{}p{0.34\linewidth}p{0.6\linewidth}@{}}
\toprule
Argumento & Efecto \\
\midrule
\endhead
\texttt{--shot ruta.png} & Captura la ventana y sale \\
\texttt{--delay ms} & Espera antes de capturar \\
\texttt{--section id} & \texttt{inicio}, \texttt{explorar}, \texttt{playlists}, \texttt{albumes}, \texttt{megusta} o \texttt{buscar} \\
\texttt{--index n} & Lleva la vista al elemento \texttt{n} \\
\texttt{--step n} (\texttt{--step-delay ms}) & Mueve la vista \texttt{n} elementos (capturar en movimiento) \\
\texttt{--search texto} & Busca al arrancar \\
\texttt{--autoplay} (\texttt{--autoplay-delay ms}) & Reproduce el elemento enfocado \\
\texttt{--then-next ms} & Pasa a la siguiente pista \\
\texttt{--open} (\texttt{--shuffle}) & Abre el detalle (y pulsa «Aleatorio») \\
\texttt{--inspect [flip|open|disc]} & Modo inspeccionar \\
\texttt{--now} & Abre la pantalla «Sonando» \\
\texttt{--tab letra|cola|temporal} & Pestaña del panel derecho \\
\texttt{--temporal} & Añade a la lista temporal y abre el menú contextual \\
\texttt{--settings}, \texttt{--settings-end} & Abre las opciones (al principio o al final) \\
\texttt{--settings-cat id} & Abre las opciones en una categoría: \texttt{cuenta}, \texttt{reproduccion}, \texttt{apariencia}, \texttt{vistas}, \texttt{plugins} \\
\texttt{--settings-at y} & Abre las opciones desplazadas \texttt{y} píxeles \\
\texttt{--minimize ms} [\texttt{--hide}|\texttt{--mini}] [\texttt{--restore ms}] & Minimiza, oculta o abre el mini; y vuelve \\
\texttt{--panel} & Con el panel de letra y cola abierto (sin él, arranca oculto) \\
\bottomrule
\end{longtable}
\end{center}

Variables de entorno útiles:

\begin{center}
\small
\begin{longtable}{@{}p{0.34\linewidth}p{0.6\linewidth}@{}}
\toprule
Variable & Efecto \\
\midrule
\endhead
\texttt{MUSICPLAYER\_SILENT=1} & mpv no saca sonido (pruebas) \\
\texttt{MUSICPLAYER\_DEBUG\_CLICK} & Clics de ratón reales: \texttt{"x,y,left|right,ms;…"}. Arrastres: \texttt{"x,y,drag,ms,x2,y2"}, o \texttt{draghold} para no soltar y capturar a mitad del gesto \\
\texttt{MUSICPLAYER\_LOGIN\_SHOT}, \texttt{\_URL} & Captura y página inicial de la ventana de inicio de sesión (macOS) \\
\texttt{MUSICPLAYER\_YTDLP\_CLIENTS} & Fuerza los clientes de YouTube de yt-dlp (experimentos) \\
\texttt{MUSICPLAYER\_DISCORD\_ID} & (al compilar) Otra aplicación de Discord \\
\texttt{HOME=/carpeta} & Aísla preferencias, sesión y caché de las tuyas \\
\texttt{HTTPS\_PROXY} & Proxy para la red \\
\texttt{QT\_QPA\_PLATFORM=offscreen} & Sin ventana visible \\
\bottomrule
\end{longtable}
\end{center}

\section{Pruebas}

\begin{terminal}
cargo test --release
\end{terminal}

Hay pruebas para lo que es lógica pura y se puede comprobar sin red: el formato de las
cookies, la firma de la sesión, la elección del mejor formato de audio, el protocolo de
Discord (contra un Discord simulado), la firma de Last.fm y sus reglas de scrobble, cómo se
convierten las tarjetas en pistas, cómo se arman los avisos de escucha para YouTube Music (y
que las cookies solo vayan a servidores de YouTube)\dots{} La interfaz se prueba con capturas y clics simulados.

\begin{trampa}[No castigues a YouTube]
Cada prueba que va a la red hace peticiones reales. Si se encadenan muchas, YouTube rechaza
la IP un rato («Requested format is not available», 403, «Sign in to confirm you're not a
bot»). Agrupa las pruebas, espera unos minutos si aparecen esos errores y no los tomes por
fallos del código sin comprobarlo antes con \texttt{yt-dlp} directamente.
\end{trampa}
